\section{Non-Functional Requirements}
\label{sec:sd-nfr}

The non-functional requirements (Table~\ref{tab:sd-nfr}) state the qualities that the needs of Section~\ref{sec:sd-stakeholders} demand from the system as a whole. They are grouped by quality characteristic, following the product quality model of ISO/IEC~25010. Most of them come from the nature of the workload: a run makes many paid LLM calls, takes minutes and waits for a human for an unbounded time, so reliability, cost control and traceability of the output weigh more than raw throughput.

The \emph{Verification} column names how each requirement is checked: by an automated test (\emph{Test}), by a static check or a review of the design and code (\emph{Inspection}), by running the system and observing it (\emph{Demonstration}), or by a study with users (\emph{User study}). For a requirement that is not met yet, it is the method planned. The \emph{Status} column follows the same convention as Table~\ref{tab:sd-fr}. NFR-05 is partly met because provider credentials are not stored yet, so their encryption is not in place; NFR-08 depends on the evaluation harness; no tracing tool is connected yet for NFR-09; NFR-10 has not been assessed with users; and NFR-15 is met by the development setup, while the production layout of Section~\ref{sec:sd-deployment} is not deployed yet.

\begingroup\setstretch{1.15}\small
\begin{longtable}{|p{0.09\linewidth}|p{0.17\linewidth}|p{0.33\linewidth}|p{0.16\linewidth}|p{0.1\linewidth}|}
    \caption{Non-functional requirements.}
    \label{tab:sd-nfr} \\
    \hline
    \textbf{ID} & \textbf{Category} & \textbf{Requirement} & \textbf{Verification} & \textbf{Status} \\
    \hline
    \endfirsthead
    \hline
    \textbf{ID} & \textbf{Category} & \textbf{Requirement} & \textbf{Verification} & \textbf{Status} \\
    \hline
    \endhead
    NFR-01 & Reliability & A closed browser does not stop a run; a reconnecting client receives every event it missed. & Test & Done \\ \hline
    NFR-02 & Reliability & After a failure, a run continues from its last saved checkpoint without repeating the LLM calls of earlier steps. & Demonstration & Done \\ \hline
    NFR-03 & Performance & The event stream sends a heartbeat every 15 seconds; streamed tokens appear in the web app as they are generated. & Demonstration & Done \\ \hline
    NFR-04 & Cost & Cost limits per run are checked before every model call; usage, including prompt cache reads and writes, is recorded for every step. & Inspection & Done \\ \hline
    NFR-05 & Security & Sessions use an httpOnly, Secure, SameSite cookie; provider keys are stored encrypted and never returned by the API; a resource without access answers 404. & Test, Inspection & Partial \\ \hline
    NFR-06 & Modifiability & Models and providers are switched by configuration, per agent role; a change of the workflow is a new versioned workflow definition with no database migration. & Inspection & Done \\ \hline
    NFR-07 & Data integrity & Every artifact and item is validated against its versioned content schema and the integrity rules before it is saved. & Test & Done \\ \hline
    NFR-08 & Reproducibility & Experiments import the same agent code as the product and record the workflow version, prompt versions, models and commit of every run. & Inspection & Partial \\ \hline
    NFR-09 & Observability & Every agent step can be inspected in the LLM tracing tool. & Demonstration & Planned \\ \hline
    NFR-10 & Usability & The current stage of a run, the pending review and the artifacts it changed are visible without leaving the thread page; a review needs no more than one screen. & User study & Not assessed \\ \hline
    NFR-11 & Maintainability & The backend is divided into modules with allowed import directions (Section~\ref{sec:sd-backend-modules}); the rule is checked automatically on every change. & Inspection & Done \\ \hline
    NFR-12 & Testability & The agent workflow can be run without the API and with simulated language models, so that every step and transition can be tested without a provider. & Test & Done \\ \hline
    NFR-13 & Interoperability & The API is described by an OpenAPI schema from which the client of the web app is generated; the requirements can be exported as Markdown, CSV and JSON. & Test & Done \\ \hline
    NFR-14 & Boundedness & Every loop of the workflow has a bound, and a run stops after at most eighty steps, so that every run terminates and its cost has an upper limit. & Inspection & Done \\ \hline
    NFR-15 & Deployability & The services run in containers orchestrated by Docker Compose, and the production baseline runs them all on one Linux host; secrets are supplied through the environment and never stored with the source. & Demonstration & Partial \\ \hline
\end{longtable}
\endgroup
